\documentclass[12pt]{article}
\usepackage{amsmath}
\usepackage{amssymb}
\usepackage{geometry}
\usepackage{amsthm}
\usepackage{enumitem}
\geometry{letterpaper, margin=1in}

\title{Homework 4.1: Vector Spaces and Subspaces}
\author{Tashfeen Omran}
\date{September 2026}

\begin{document}

\maketitle
\tableofcontents
\newpage

\section{Part 1: Comprehensive Introduction, Context, and Prerequisites}

\subsection{Core Concepts}
In previous chapters, we studied vectors in $\mathbb{R}^n$ (lists of $n$ real numbers) and the operations of vector addition and scalar multiplication. The topic of \textbf{Vector Spaces} takes these familiar operations and abstracts them. A vector space is simply \textit{any} collection of objects (which we call "vectors") for which we can define addition and scalar multiplication such that they satisfy a specific set of ten rules (axioms). 

These "vectors" do not have to be columns of numbers. They can be polynomials, continuous functions, matrices, or even infinite sequences. If a set of objects and their operations satisfy the ten axioms, all the algebraic theorems we proved for $\mathbb{R}^n$ automatically apply to them.

A \textbf{Subspace} is a smaller vector space that lives entirely inside a larger vector space. Because the larger space already guarantees that most of the ten axioms hold (like commutativity and associativity), checking if a subset is a subspace only requires verifying three specific properties: it must contain the zero vector, it must be closed under addition, and it must be closed under scalar multiplication.

\subsection{Intuition and Derivation}
The intuition behind a vector space is the concept of "linearity." If you have two signals (like audio waves), you can add them together to get a new signal, or you can amplify one by multiplying it by a constant. The set of all possible signals forms a vector space because these operations behave exactly like adding or scaling arrows in 2D space. 

The concept of a subspace derives from the idea of a "flat" geometric object passing through the origin. In $\mathbb{R}^3$, a plane passing through the origin is a subspace. If you add any two vectors on that plane, the result is still on the plane (closure under addition). If you stretch a vector on that plane, it remains on the plane (closure under scalar multiplication). The origin itself must be included (the zero vector).

\subsection{Historical Context and Motivation}
The formal definition of an abstract vector space was first given by the Italian mathematician Giuseppe Peano in 1888. Before this, mathematicians worked with specific types of vectors (like geometric arrows or $n$-tuples). However, as the study of differential equations and function analysis advanced in the 19th century, mathematicians realized that sets of functions behaved algebraically just like $\mathbb{R}^n$. Peano's axiomatization allowed mathematicians to unify the study of geometry, algebra, and calculus. By proving a theorem once for an abstract vector space, it was simultaneously proven for matrices, polynomials, and solutions to differential equations.

\subsection{Key Formulas and Definitions}
\textbf{The Subspace Test:} A subset $H$ of a vector space $V$ is a subspace if and only if it satisfies three conditions:
\begin{enumerate}
    \item The zero vector of $V$ is in $H$ ($0 \in H$).
    \item $H$ is closed under vector addition: For each $u$ and $v$ in $H$, the sum $u + v$ is in $H$.
    \item $H$ is closed under scalar multiplication: For each $u$ in $H$ and each scalar $c$, the vector $cu$ is in $H$.
\end{enumerate}

\textbf{Span as a Subspace:} If $v_1, \dots, v_p$ are in a vector space $V$, then $\text{Span}\{v_1, \dots, v_p\}$ is a subspace of $V$.

\subsection{Prerequisites}
To succeed in this chapter, you must be proficient in:
\begin{itemize}
    \item Solving systems of linear equations using augmented matrices and row reduction (RREF).
    \item Understanding the definition of the $\text{Span}$ of a set of vectors.
    \item Basic polynomial algebra and matrix addition/scalar multiplication.
\end{itemize}

\newpage

\section{Part 2: Detailed Homework Solutions}
\textit{Note: A strategic selection of problems has been chosen to cover every distinct problem archetype found in the exercises, ensuring comprehensive mastery without redundant repetition.}

\subsection{Problem 1}
\textbf{Statement:} Let $V$ be the first quadrant in the $xy$-plane; that is, let $V = \left\{ \begin{bmatrix} x \\ y \end{bmatrix} : x \ge 0, y \ge 0 \right\}$. If $u$ and $v$ are in $V$, is $u+v$ in $V$? Why? Find a specific vector $u$ in $V$ and a specific scalar $c$ such that $cu$ is not in $V$. (This is enough to show that $V$ is not a vector space.)

\textbf{Solution:}
First, we check closure under addition. Let $u = \begin{bmatrix} x_1 \\ y_1 \end{bmatrix}$ and $v = \begin{bmatrix} x_2 \\ y_2 \end{bmatrix}$ be in $V$. This means $x_1, y_1, x_2, y_2 \ge 0$.
The sum is $u + v = \begin{bmatrix} x_1 + x_2 \\ y_1 + y_2 \end{bmatrix}$.
Since the sum of two non-negative numbers is non-negative, $x_1 + x_2 \ge 0$ and $y_1 + y_2 \ge 0$. Thus, $u+v \in V$. Yes, it is closed under addition.

Next, we find a counterexample for scalar multiplication. 
Let $u = \begin{bmatrix} 1 \\ 1 \end{bmatrix}$. Since $1 \ge 0$ and $1 \ge 0$, $u \in V$.
Let the scalar $c = -1$.
Then $cu = -1 \begin{bmatrix} 1 \\ 1 \end{bmatrix} = \begin{bmatrix} -1 \\ -1 \end{bmatrix}$.
Since the entries are negative, $cu \notin V$. 
Because $V$ is not closed under scalar multiplication, it is \textbf{not a vector space}.

\subsection{Problem 2}
\textbf{Statement:} Let $W$ be the union of the first and third quadrants in the $xy$-plane. That is, let $W = \left\{ \begin{bmatrix} x \\ y \end{bmatrix} : xy \ge 0 \right\}$. 
a. If $u$ is in $W$ and $c$ is any scalar, is $cu$ in $W$? Why?
b. Find specific vectors $u$ and $v$ in $W$ such that $u+v$ is not in $W$.

\textbf{Solution:}
\textbf{Part a:} Let $u = \begin{bmatrix} x \\ y \end{bmatrix} \in W$, meaning $xy \ge 0$. Let $c$ be any real scalar.
$cu = \begin{bmatrix} cx \\ cy \end{bmatrix}$.
To check if $cu \in W$, we multiply its components: $(cx)(cy) = c^2(xy)$.
Since $c^2 \ge 0$ for all real $c$, and $xy \ge 0$ by assumption, their product $c^2(xy) \ge 0$. 
Thus, $cu \in W$. Yes, $W$ is closed under scalar multiplication.

\textbf{Part b:} We need a counterexample for addition. We must pick one vector from the first quadrant and one from the third quadrant such that their sum lands in the second or fourth quadrant.
Let $u = \begin{bmatrix} 1 \\ 4 \end{bmatrix} \in W$ (since $1 \cdot 4 = 4 \ge 0$).
Let $v = \begin{bmatrix} -4 \\ -1 \end{bmatrix} \in W$ (since $(-4) \cdot (-1) = 4 \ge 0$).
Now, compute $u + v$:
\[ u + v = \begin{bmatrix} 1 \\ 4 \end{bmatrix} + \begin{bmatrix} -4 \\ -1 \end{bmatrix} = \begin{bmatrix} -3 \\ 3 \end{bmatrix} \]
Check if $u+v \in W$ by multiplying its components: $(-3)(3) = -9$.
Since $-9 < 0$, $u+v \notin W$. Therefore, $W$ is \textbf{not a vector space}.

\subsection{Problem 5}
\textbf{Statement:} Determine if the given set is a subspace of $\mathbb{P}_n$ for an appropriate value of $n$. Justify your answer.
All polynomials of the form $p(t) = at^2$, where $a$ is in $\mathbb{R}$.

\textbf{Solution:}
Let $H$ be the set of all polynomials of the form $p(t) = at^2$. We apply the three-part subspace test.
1. \textbf{Zero Vector:} The zero vector in $\mathbb{P}_n$ is the zero polynomial, $0(t) = 0$. If we choose $a = 0$, then $p(t) = 0t^2 = 0$. Thus, the zero polynomial is in $H$.
2. \textbf{Closure under Addition:} Let $p(t) = at^2$ and $q(t) = bt^2$ be two polynomials in $H$.
\[ p(t) + q(t) = at^2 + bt^2 = (a+b)t^2 \]
Since $a+b$ is a real number, the sum is of the form $(\text{constant})t^2$, so $p(t) + q(t) \in H$.
3. \textbf{Closure under Scalar Multiplication:} Let $p(t) = at^2 \in H$ and let $c$ be a real scalar.
\[ c \cdot p(t) = c(at^2) = (ca)t^2 \]
Since $ca$ is a real number, the scalar multiple is of the form $(\text{constant})t^2$, so $c \cdot p(t) \in H$.

Since all three conditions are met, the set is a \textbf{subspace} of $\mathbb{P}_2$ (and any $\mathbb{P}_n$ for $n \ge 2$).

\subsection{Problem 6}
\textbf{Statement:} Determine if the given set is a subspace of $\mathbb{P}_n$.
All polynomials of the form $p(t) = a + t^2$, where $a$ is in $\mathbb{R}$.

\textbf{Solution:}
Let $H$ be the set of polynomials $p(t) = a + t^2$. We check the first condition of the subspace test: does $H$ contain the zero vector?
The zero vector is the polynomial $0(t) = 0$ for all $t$.
For $p(t)$ to be the zero polynomial, we must have $a + t^2 = 0$ for all $t$.
This implies $a = -t^2$. However, $a$ must be a constant real number, and $-t^2$ is a variable expression. There is no constant $a$ that makes $a + t^2 = 0$ for all $t$.
Therefore, the zero polynomial is not in $H$. 
Because it fails the zero vector test, the set is \textbf{not a subspace}.

\subsection{Problem 9}
\textbf{Statement:} Let $H$ be the set of all vectors of the form $\begin{bmatrix} 3s \\ 2s \end{bmatrix}$. Find a vector $v$ in $\mathbb{R}^2$ such that $H = \text{Span}\{v\}$. Why does this show that $H$ is a subspace of $\mathbb{R}^2$?

\textbf{Solution:}
We can factor out the scalar $s$ from the vector definition:
\[ \begin{bmatrix} 3s \\ 2s \end{bmatrix} = s \begin{bmatrix} 3 \\ 2 \end{bmatrix} \]
This shows that every vector in $H$ is a scalar multiple of the vector $v = \begin{bmatrix} 3 \\ 2 \end{bmatrix}$.
Therefore, $H = \text{Span}\left\{ \begin{bmatrix} 3 \\ 2 \end{bmatrix} \right\}$.
By Theorem 1 in Section 4.1, the span of any set of vectors in a vector space $V$ is automatically a subspace of $V$. Since $H$ is the span of a vector in $\mathbb{R}^2$, $H$ is a \textbf{subspace of $\mathbb{R}^2$}.

\subsection{Problem 13}
\textbf{Statement:} Let $v_1 = \begin{bmatrix} 1 \\ -1 \\ -2 \end{bmatrix}$, $v_2 = \begin{bmatrix} 5 \\ -4 \\ -7 \end{bmatrix}$, $v_3 = \begin{bmatrix} -3 \\ 1 \\ 0 \end{bmatrix}$, and $y = \begin{bmatrix} -4 \\ 3 \\ h \end{bmatrix}$. For what value(s) of $h$ will $y$ be in the subspace of $\mathbb{R}^3$ spanned by $v_1, v_2, v_3$?

\textbf{Solution:}
The vector $y$ is in $\text{Span}\{v_1, v_2, v_3\}$ if and only if there exist scalars $c_1, c_2, c_3$ such that:
\[ c_1v_1 + c_2v_2 + c_3v_3 = y \]
This is equivalent to solving the linear system with the augmented matrix $[v_1 \ v_2 \ v_3 \ | \ y]$.
\[ \begin{bmatrix} 1 & 5 & -3 & -4 \\ -1 & -4 & 1 & 3 \\ -2 & -7 & 0 & h \end{bmatrix} \]
We row reduce this matrix to echelon form.
$R_2 \leftarrow R_2 + R_1$:
\[ \begin{bmatrix} 1 & 5 & -3 & -4 \\ 0 & 1 & -2 & -1 \\ -2 & -7 & 0 & h \end{bmatrix} \]
$R_3 \leftarrow R_3 + 2R_1$:
\[ \begin{bmatrix} 1 & 5 & -3 & -4 \\ 0 & 1 & -2 & -1 \\ 0 & 3 & -6 & h-8 \end{bmatrix} \]
$R_3 \leftarrow R_3 - 3R_2$:
\[ \begin{bmatrix} 1 & 5 & -3 & -4 \\ 0 & 1 & -2 & -1 \\ 0 & 0 & 0 & h-8 - 3(-1) \end{bmatrix} = \begin{bmatrix} 1 & 5 & -3 & -4 \\ 0 & 1 & -2 & -1 \\ 0 & 0 & 0 & h-5 \end{bmatrix} \]
For the system to be consistent (meaning a solution exists and $y$ is in the span), the last row cannot be of the form $[0 \ 0 \ 0 \ | \ b]$ where $b \neq 0$. 
Therefore, we must have:
\[ h - 5 = 0 \implies h = 5 \]
Thus, $y$ is in the subspace spanned by $v_1, v_2, v_3$ \textbf{if and only if $h = 5$}.

\subsection{Problem 21}
\textbf{Statement:} Determine if the set $H$ of all matrices of the form $\begin{bmatrix} a & b \\ 0 & d \end{bmatrix}$ is a subspace of $M_{2 \times 2}$.

\textbf{Solution:}
We apply the three-part subspace test for $H$.
1. \textbf{Zero Vector:} The zero vector in $M_{2 \times 2}$ is the $2 \times 2$ zero matrix. If we let $a=0, b=0, d=0$, we get $\begin{bmatrix} 0 & 0 \\ 0 & 0 \end{bmatrix}$, which is in $H$.
2. \textbf{Closure under Addition:} Let $A = \begin{bmatrix} a_1 & b_1 \\ 0 & d_1 \end{bmatrix}$ and $B = \begin{bmatrix} a_2 & b_2 \\ 0 & d_2 \end{bmatrix}$ be in $H$.
\[ A + B = \begin{bmatrix} a_1+a_2 & b_1+b_2 \\ 0 & d_1+d_2 \end{bmatrix} \]
Since the bottom-left entry is $0$, the resulting matrix is of the required form. Thus, $A+B \in H$.
3. \textbf{Closure under Scalar Multiplication:} Let $A = \begin{bmatrix} a & b \\ 0 & d \end{bmatrix} \in H$ and let $c$ be a scalar.
\[ cA = \begin{bmatrix} ca & cb \\ 0 & cd \end{bmatrix} \]
The bottom-left entry remains $0$, so $cA \in H$.

Since all three conditions are satisfied, $H$ is a \textbf{subspace} of $M_{2 \times 2}$.

\subsection{Problem 33}
\textbf{Statement:} Complete the following proof that the zero vector is unique. Suppose that $w$ in $V$ has the property that $u + w = w + u = u$ for all $u$ in $V$. In particular, $0 + w = 0$. But $0 + w = w$, by Axiom \_\_\_\_\_. Hence $w = 0 + w = 0$.

\textbf{Solution:}
The problem provides the structure of the proof. We are given a vector $w$ that acts like a zero vector (i.e., $u+w=u$ for any $u$). We want to prove $w$ is exactly the standard zero vector $0$.
Step 1: Since $w$ acts like a zero vector for any $u$, we can choose $u = 0$. This gives $0 + w = 0$.
Step 2: By the standard vector space Axiom 4, the standard zero vector $0$ has the property that $0 + v = v$ for any vector $v$. Applying this to $v = w$, we get $0 + w = w$.
Step 3: Equating the two results for $0 + w$, we get $w = 0$.
The missing axiom in the blank is \textbf{Axiom 4} (the existence of a zero vector).

\newpage

\section{Part 3: In-Depth Analysis of Problems and Techniques}

\subsection{A) Problem Types and General Approach}
\begin{itemize}
    \item \textbf{Type 1: Counterexample Generation (Problems 1, 2, 6)}
    \textit{Strategy:} When a set is restricted by inequalities (like $x \ge 0$) or non-linear conditions (like $xy \ge 0$ or $+ t^2$), it is highly likely it is NOT a subspace. The fastest way to prove this is to find a single numerical counterexample where addition or scalar multiplication forces the vector outside the set. Checking the zero vector (Problem 6) is always the easiest first step.
    
    \item \textbf{Type 2: Subspace Verification via Axioms (Problems 5, 21)}
    \textit{Strategy:} When a set is defined by linear combinations of variables, it is usually a subspace. To prove it, you must rigorously apply the 3-step subspace test. Define two generic elements (using subscripts like $a_1, a_2$), add them, and show the result matches the structural definition of the set. Repeat for scalar multiplication.
    
    \item \textbf{Type 3: Span Identification (Problem 9)}
    \textit{Strategy:} If a set is defined by a vector with parameters (like $s$ and $t$), factor the parameters out to express the set as a linear combination of constant vectors. Once written as a $\text{Span}$, you can immediately invoke Theorem 1 to declare it a subspace without doing the 3-step test.
    
    \item \textbf{Type 4: Span Membership via Linear Systems (Problem 13)}
    \textit{Strategy:} To check if a vector $y$ is in a subspace spanned by $\{v_1, \dots, v_p\}$, set up the augmented matrix $[v_1 \dots v_p \ | \ y]$. Row reduce to echelon form. If the system is consistent (no row of $[0 \dots 0 \ | \ b]$ with $b \neq 0$), the vector is in the subspace.
\end{itemize}

\subsection{B) Key Algebraic and Calculus Manipulations}
\begin{itemize}
    \item \textbf{Parameter Factoring:} In Problem 9, the algebraic trick of factoring a scalar out of a vector $\begin{bmatrix} 3s \\ 2s \end{bmatrix} = s \begin{bmatrix} 3 \\ 2 \end{bmatrix}$ is crucial. It translates a set description into a formal $\text{Span}$, bridging the gap between algebra and vector space theory.
    \item \textbf{Row Reduction for Consistency:} In Problem 13, we didn't need to find the exact values of $c_1, c_2, c_3$. We only needed to know IF a solution exists. Stopping at echelon form (rather than reduced echelon form) is a time-saving algebraic technique. We set the expression in the augmented column ($h-5$) to $0$ to force consistency.
    \item \textbf{Polynomial Coefficient Matching:} In Problem 6, setting $a + t^2 = 0$ and realizing that a constant $a$ cannot equal a variable $-t^2$ relies on the algebraic principle that two polynomials are equal if and only if all their corresponding coefficients are equal.
\end{itemize}

\newpage

\section{Part 4: "Cheatsheet" and Tips for Success}

\subsection{Summary of Important Formulas and Theorems}
\begin{itemize}
    \item \textbf{The 3-Step Subspace Test:} To prove $H$ is a subspace of $V$:
        1. Show $0 \in H$.
        2. Show $u, v \in H \implies u+v \in H$.
        3. Show $u \in H, c \in \mathbb{R} \implies cu \in H$.
    \item \textbf{Theorem 1:} If $v_1, \dots, v_p \in V$, then $\text{Span}\{v_1, \dots, v_p\}$ is a subspace of $V$.
    \item \textbf{Standard Vector Spaces:} $\mathbb{R}^n$ (column vectors), $\mathbb{P}_n$ (polynomials of degree $\le n$), $M_{m \times n}$ (matrices), $C[a,b]$ (continuous functions on an interval).
\end{itemize}

\subsection{Tricks and Shortcuts}
\begin{itemize}
    \item \textbf{The Zero Vector Cheat:} Always check if the zero vector is in the set first. If it isn't, you are done; it's not a subspace. This saves you from doing the addition and scalar multiplication proofs.
    \item \textbf{The Span Shortcut:} If you can write a set as the span of some vectors, you don't need to do the 3-step test. Just cite Theorem 1.
    \item \textbf{Linearity Indicator:} If the definition of a set involves squares ($x^2$), absolute values ($|x|$), inequalities ($\ge$), or constants added to variables ($x+1$), it is almost certainly \textbf{not} a subspace.
\end{itemize}

\subsection{Common Pitfalls}
\begin{itemize}
    \item \textbf{Confusing "Vector" with "Array":} Remember that a polynomial $p(t) = t^2$ or a matrix $A$ is considered a "vector" in this chapter. Do not try to force them into column brackets.
    \item \textbf{Proving by Example:} You cannot prove a set IS a subspace by picking specific numbers (e.g., "Let $u=1, v=2$, it works!"). You must use generic variables. However, you \textit{can} (and should) use specific numbers to prove a set is NOT a subspace (counterexample).
    \item \textbf{Forgetting the Parent Space:} A subspace must be a subset of a known vector space. $\mathbb{R}^2$ is not a subspace of $\mathbb{R}^3$ because a 2D vector $\begin{bmatrix} x \\ y \end{bmatrix}$ is not a 3D vector. However, the set of vectors $\begin{bmatrix} x \\ y \\ 0 \end{bmatrix}$ IS a subspace of $\mathbb{R}^3$.
\end{itemize}

\newpage

\section{Part 5: Conceptual Synthesis and The "Big Picture"}

\subsection{A) Thematic Connections}
The core theme of this topic is \textbf{Abstraction and Generalization}. In mathematics, once we understand a concrete system (like arrows in $\mathbb{R}^2$), we extract the rules that govern that system (the 10 axioms) and look for other systems that follow the exact same rules. 

This theme appears everywhere in calculus. For example, when we learned the derivative, we first applied it to simple polynomials. Then we abstracted the rules of differentiation (linearity, product rule, chain rule) and applied them to trigonometric and exponential functions. Vector spaces do the same thing for algebra: they take the rules of $\mathbb{R}^n$ and apply them to functions and matrices.

\subsection{B) Forward and Backward Links}
\textbf{Backward Links:} This chapter relies heavily on Chapter 1 (Linear Combinations and Spans). The concept of a subspace is essentially a geometric realization of a span. When we solved $Ax = 0$ in Chapter 1, we were actually finding a subspace (which we will soon call the Null Space).

\textbf{Forward Links:} Understanding abstract vector spaces is the absolute bedrock for the rest of linear algebra and advanced mathematics. 
\begin{itemize}
    \item \textbf{Eigenvectors:} Later, we will find "eigenspaces," which are subspaces of vectors that only scale when multiplied by a matrix.
    \item \textbf{Differential Equations:} The set of all solutions to a homogeneous linear differential equation (like $y'' + y = 0$) forms a vector space. The "fundamental set of solutions" you find in a differential equations class is actually just a \textit{basis} for this vector space.
    \item \textbf{Fourier Series:} In advanced calculus, we treat functions like $\sin(nx)$ and $\cos(nx)$ as vectors in an infinite-dimensional vector space, allowing us to project complex waves onto these basic "vectors."
\end{itemize}

\newpage

\section{Part 6: Abstract Real-World Modeling (AI, ML, \& Quant Finance)}

\subsection{A) Concrete Scenarios}
\begin{enumerate}
    \item \textbf{Support Vector Machines (SVM) and Reproducing Kernel Hilbert Spaces (ML):} In machine learning, when data is not linearly separable in its original space, we map it to a higher-dimensional abstract vector space (a Hilbert space) where it becomes separable. The mathematics of these spaces relies entirely on the axioms of vector spaces and inner products.
    \item \textbf{Portfolio Optimization (Quant Finance):} A portfolio can be represented as a vector of weights $w \in \mathbb{R}^n$. The set of all self-financing portfolios (where weights sum to 0, representing long/short positions) forms a subspace. Quantitative analysts project risk vectors onto this subspace to hedge against market movements.
    \item \textbf{State Spaces in Reinforcement Learning (AI):} In continuous-space Markov Decision Processes, the value function $V(s)$ is an element of a vector space of functions. Bellman equations act as linear operators on this vector space, and finding the optimal policy is equivalent to finding the fixed point of this operator in the function space.
\end{enumerate}

\subsection{B) Model Problem Setup: RKHS for Support Vector Machines}
Let us rigorously model the feature mapping in an SVM using vector space theory.

Let $\mathcal{X} \subseteq \mathbb{R}^d$ be our input space (e.g., images). We define a feature map $\phi: \mathcal{X} \to \mathcal{H}$, where $\mathcal{H}$ is a high-dimensional (possibly infinite) abstract vector space called a Hilbert space. Because $\mathcal{H}$ is a vector space, it satisfies closure under addition and scalar multiplication.

We define an inner product on this vector space, $\langle \cdot, \cdot \rangle_{\mathcal{H}}$. The \textit{Kernel trick} allows us to compute inner products in $\mathcal{H}$ without explicitly mapping the data, using a kernel function $K(x, y) = \langle \phi(x), \phi(y) \rangle_{\mathcal{H}}$.

The SVM optimization problem seeks a hyperplane in the vector space $\mathcal{H}$ that separates two classes of data with maximum margin. The normal vector to this hyperplane, $w$, is an element of the vector space $\mathcal{H}$. 
The primal optimization problem is formulated as:
\[ \min_{w \in \mathcal{H}, b \in \mathbb{R}} \frac{1}{2} \|w\|_{\mathcal{H}}^2 \]
subject to the constraints:
\[ y_i (\langle w, \phi(x_i) \rangle_{\mathcal{H}} + b) \ge 1 \quad \text{for all } i = 1, \dots, N \]
By the Representer Theorem, the optimal vector $w^*$ lies in the \textbf{subspace} spanned by the mapped training data:
\[ w^* \in \text{Span}\{ \phi(x_1), \dots, \phi(x_N) \} \subset \mathcal{H} \]
Thus, $w^* = \sum_{i=1}^N \alpha_i y_i \phi(x_i)$. This demonstrates how the concept of a subspace and span directly dictates the architecture of one of the most powerful classification algorithms in machine learning.

\newpage

\section{Part 7: Common Variations and Untested Concepts}

\subsection{Intersection of Two Subspaces}
A very common theoretical question not explicitly solved in our homework selection is proving that the intersection of two subspaces is always a subspace.

\textbf{Theorem:} If $H$ and $K$ are subspaces of a vector space $V$, then their intersection $H \cap K$ is also a subspace of $V$.

\textbf{Proof Example:}
Let $H$ and $K$ be subspaces of $V$. We apply the 3-step test to $H \cap K$.
1. \textbf{Zero Vector:} Since $H$ is a subspace, $0 \in H$. Since $K$ is a subspace, $0 \in K$. Therefore, $0 \in H \cap K$.
2. \textbf{Addition:} Let $u, v \in H \cap K$. This means $u, v \in H$ and $u, v \in K$. Because $H$ is a subspace, $u+v \in H$. Because $K$ is a subspace, $u+v \in K$. Therefore, $u+v \in H \cap K$.
3. \textbf{Scalar Multiplication:} Let $u \in H \cap K$ and $c \in \mathbb{R}$. Then $u \in H$ and $u \in K$. Because $H$ and $K$ are subspaces, $cu \in H$ and $cu \in K$. Therefore, $cu \in H \cap K$.
Since all three conditions hold, $H \cap K$ is a subspace.

\subsection{Union of Two Subspaces}
Conversely, the union of two subspaces $H \cup K$ is \textbf{rarely} a subspace. It is only a subspace if one is entirely contained within the other. 
\textbf{Example:} Let $H$ be the $x$-axis in $\mathbb{R}^2$ and $K$ be the $y$-axis. Both are subspaces. Their union is the set of all vectors on the axes. Let $u = \begin{bmatrix} 1 \\ 0 \end{bmatrix} \in H$ and $v = \begin{bmatrix} 0 \\ 1 \end{bmatrix} \in K$. Both are in $H \cup K$. But $u+v = \begin{bmatrix} 1 \\ 1 \end{bmatrix}$, which is on neither axis, so it is not in $H \cup K$. Thus, the union fails closure under addition.

\newpage

\section{Part 8: Advanced Diagnostic Testing: "Find the Flaw"}
\textit{Instructions: Each of the following solutions contains a single, subtle conceptual or algebraic flaw. Identify the error, explain why it is wrong, and provide the correct solution.}

\subsection{Problem 1}
\textbf{Question:} Let $H = \left\{ \begin{bmatrix} x \\ y \end{bmatrix} : x + y = 1 \right\}$. Is $H$ a subspace of $\mathbb{R}^2$?
\textbf{Flawed Solution:}
We check closure under scalar multiplication. Let $u = \begin{bmatrix} x \\ y \end{bmatrix} \in H$, so $x+y=1$. Let $c=1$.
$cu = 1 \begin{bmatrix} x \\ y \end{bmatrix} = \begin{bmatrix} x \\ y \end{bmatrix}$.
The sum of the components is $x+y = 1$. Since the sum is 1, $cu \in H$. 
Therefore, $H$ is closed under scalar multiplication and is a subspace.

\textbf{Correction:}
\textbf{Flaw:} The student chose a specific scalar $c=1$ to prove closure, but closure under scalar multiplication must hold for \textit{all} real scalars $c$.
\textbf{Correct Solution:} $H$ is not a subspace. The easiest way to show this is that it does not contain the zero vector. For the zero vector $\begin{bmatrix} 0 \\ 0 \end{bmatrix}$, $x+y = 0+0 = 0 \neq 1$. Thus $0 \notin H$.

\subsection{Problem 2}
\textbf{Question:} Is $y = \begin{bmatrix} 1 \\ 2 \end{bmatrix}$ in the span of $v_1 = \begin{bmatrix} 2 \\ 4 \end{bmatrix}$ and $v_2 = \begin{bmatrix} -1 \\ -2 \end{bmatrix}$?
\textbf{Flawed Solution:}
We set up the augmented matrix and row reduce:
\[ \begin{bmatrix} 2 & -1 & 1 \\ 4 & -2 & 2 \end{bmatrix} \xrightarrow{R_2 - 2R_1} \begin{bmatrix} 2 & -1 & 1 \\ 0 & 0 & 0 \end{bmatrix} \]
Because we got a row of all zeros, the system does not have a unique solution. Therefore, $y$ is not in the span of $v_1$ and $v_2$.

\textbf{Correction:}
\textbf{Flaw:} A row of all zeros means the system has infinitely many solutions, which means it \textit{is} consistent, not inconsistent.
\textbf{Correct Solution:} The system is consistent because there is no pivot in the rightmost column (no row of the form $[0 \ 0 \ | \ b]$ with $b \neq 0$). Since the system is consistent, $y$ \textbf{is} in the span of $v_1$ and $v_2$.

\subsection{Problem 3}
\textbf{Question:} Let $H$ be the set of all polynomials in $\mathbb{P}_2$ such that $p(0) = 0$. Is $H$ a subspace?
\textbf{Flawed Solution:}
To be a subspace, $H$ must contain the zero vector. The zero vector in $\mathbb{P}_2$ is the number $0$. But $H$ is a set of polynomials, and $0$ is a number, not a polynomial. Therefore, $H$ does not contain the zero vector and is not a subspace.

\textbf{Correction:}
\textbf{Flaw:} The student confused the scalar number $0$ with the zero polynomial $p(t) = 0$.
\textbf{Correct Solution:} The zero vector in $\mathbb{P}_2$ is the zero polynomial $p(t) = 0$ for all $t$. We check if it satisfies the condition: $p(0) = 0$. Yes, it does. Furthermore, if $p(0)=0$ and $q(0)=0$, then $(p+q)(0) = p(0)+q(0) = 0$, and $(cp)(0) = c(p(0)) = 0$. Thus, $H$ \textbf{is} a subspace.

\subsection{Problem 4}
\textbf{Question:} Let $H$ be the set of all invertible $2 \times 2$ matrices. Is $H$ a subspace of $M_{2 \times 2}$?
\textbf{Flawed Solution:}
We check closure. Let $A$ and $B$ be invertible matrices in $H$. We know from Chapter 2 that the product of two invertible matrices is invertible (i.e., $(AB)^{-1} = B^{-1}A^{-1}$). Therefore, the combination of $A$ and $B$ is in $H$. Thus, $H$ is a subspace.

\textbf{Correction:}
\textbf{Flaw:} The student checked closure under matrix \textit{multiplication}, but vector spaces require closure under matrix \textit{addition}.
\textbf{Correct Solution:} $H$ is not a subspace. First, the zero matrix is not invertible, so $0 \notin H$. Second, it is not closed under addition. Let $A = I$ and $B = -I$. Both are invertible, but $A+B = 0$, which is not invertible.

\subsection{Problem 5}
\textbf{Question:} Let $V$ be the vector space of all real-valued functions. Let $H = \{ f \in V : f(1) = 1 \}$. Is $H$ a subspace?
\textbf{Flawed Solution:}
We check scalar multiplication. Let $f \in H$, so $f(1) = 1$. Let $c$ be a scalar.
We evaluate $c \cdot f(1) = c(1) = c$.
Since $c$ is just a real number, the function is still a real-valued function. Therefore, it is closed under scalar multiplication and is a subspace.

\textbf{Correction:}
\textbf{Flaw:} The student misunderstood the condition for being in $H$. To be in $H$, the new function $(cf)$ must evaluate to $1$ at $t=1$, not just be a real number.
\textbf{Correct Solution:} $H$ is not a subspace. Let $f \in H$ (so $f(1)=1$) and let $c=2$. The new function is $g = 2f$. We check if $g \in H$ by evaluating at $t=1$: $g(1) = 2f(1) = 2(1) = 2$. Since $g(1) \neq 1$, $g \notin H$. It fails closure under scalar multiplication. (It also fails the zero vector test, as the zero function $z(t)=0$ gives $z(1)=0 \neq 1$).

\end{document}